%%%%%%%%%%%%%%%%%%%%%%%%%%%%%%%%%%%%%%%%%%%%%%%%%%%%%%%%%%%%%%%%%%%%%%
% Lớp 10 | Hàm số bậc hai và đồ thị | Hàm số và đồ thị
% Dạng 0D3?1-2: Tập xác định của hàm số
% Số câu: 48  (sinh tự động bởi tach_cau.py; từ nay thầy sửa trực tiếp file này)
%%%%%%%%%%%%%%%%%%%%%%%%%%%%%%%%%%%%%%%%%%%%%%%%%%%%%%%%%%%%%%%%%%%%%%

%%%%%%%%%%%%%%%% TRẮC NGHIỆM 4 PHƯƠNG ÁN %%%%%%%%%%%%%%%%

% nguồn: MaHoaNH_TracNghiem_Lop10
\begin{ex}%[0D3N1-2]%[ma:7e67ca5d]
	Tìm tập xác định  của hàm số $y=\dfrac{x-2}{x-1}$.
	 \choice
		 {\True $\mathscr{D}=\mathbb{R} \setminus \{1\}$}
		 {$\mathscr{D}=\mathbb{R} \setminus \{1;2\}$}
		 {$\mathscr{D}=\mathbb{R}$}
		{$\mathscr{D}=\mathbb{R} \setminus  \{2\}$}
	 \loigiai{
		Điều kiện xác định của hàm số là  $x-1\ne 0\Leftrightarrow x\ne 1$.\\
	Vậy  $\mathscr{D}=\mathbb{R} \setminus\{ 1\}$.}
\end{ex}

% nguồn: MaHoaNH_TracNghiem_Lop10
\begin{ex}%[0D3N1-2]%[ma:93027f25]
	Tìm tập xác định  $D$ của hàm số $y=\sqrt{x+2}-\sqrt{x+3}$.
	 \choice
		{$\mathscr{D}=[-3;+\infty)$} 
		 {\True $\mathscr{D}=[-2;+\infty)$}
		 {$\mathscr{D}=\mathbb{R}$}
		 {$\mathscr{D}=[2;+\infty)$}
	 \loigiai{
		Hàm số xác định khi và  chỉ khi $\heva{&x+2\ge 0\\&x+3\ge 0} \Leftrightarrow x\ge-2.$\\
	 Vậy $\mathscr{D}=[-2;+\infty)$.}
\end{ex}

% nguồn: MaHoaNH_TracNghiem_Lop10
\begin{ex}%[0D3N1-2]%[ma:e5208705]
	Cặp số nào sau đây là nghiệm của hệ  bất phương trình $\heva{& x+2y\le8\\ & 3x-y>3 }$ ?
	\choice 
		 {$\left(0; 4\right)$ }
		 {$\left(1; 1\right)$}
		 {\True $\left(4; 1\right)$}
		{$\left(2;  3\right)$}
	 \loigiai{
		Lần lượt thay các  bộ số vào hệ bất phương trình ta được một nghiệm của hệ bất phương trình trên là $\left(4; 1\right)$.
	} 
\end{ex}

% nguồn: MaHoaNH_TracNghiem_Lop10
\begin{ex}%[0D3N1-2]%[ma:c113ec07]
	Tập  xác định của hàm số $f(x)=\sqrt{-x^{2}+5x+6}$.
	\choice 
		 {$\mathbb{R}$}
		 {\True $\left[-1;6\right]$}
		 {$\left(-1;6\right)$}
		{$\left(-\infty;-1\right]\cup\left[6;+\infty\right)$} 
	 \loigiai{
		Hàm số xác định $\Leftrightarrow -x^{2}+5x+6\ge  0 \Leftrightarrow -1\le x\le 6$.
	} 
\end{ex}

% nguồn: MaHoaNH_TracNghiem_Lop10
\begin{ex}%[0D3N1-2]%[ma:ec64642b]
	Tìm tập  xác định $\mathscr{D}$ của hàm số $y=\dfrac{x}{x+2}$.
	 \choice
		{\True $\mathscr{D}=\mathbb{R} \setminus \{-2\}$} 
		{$\mathscr{D}=(-2;+\infty)$} 
		{$\mathscr{D}=\mathbb{R}  \setminus \{2\}$}
		{$\mathscr{D}=(-\infty;-2)$} 
	\loigiai{ 
		Điều kiện xác định  của hàm số là $x+2\ne 0\Leftrightarrow x\ne-2$.\\
	Vậy tập xác  định của hàm số là $\mathscr{D}=\mathbb{R} \setminus \{-2\}$.}
\end{ex}

% nguồn: MaHoaNH_TracNghiem_Lop10
\begin{ex}%[0D3N1-2]%[ma:2a9bfbb1]
	Miền nghiệm của hệ bất phương trình $\heva{& x-2y+2>0\\ & 2x+y>3 }$ là phần  mặt phẳng chứa điểm
	 \choice
		{$\left(1;1\right)$ } 
		 {$\left(-1;2\right)$}
		{$\left(2;-1\right)$} 
		{\True $\left(2;1 \right)$} 
	 \loigiai{
		Lần lượt thay các bộ số vào hệ bất phương trình ta được một nghiệm của hệ bất phương trình trên là $\left(2;1\right)$. 
	 }
\end{ex}

% nguồn: MaHoaNH_TracNghiem_Lop10
\begin{ex}%[0D3N1-2]%[ma:564db08e]
	Tìm  tập xác định $D$ của hàm số $y=\dfrac{3x-1}{2x^2-3x-5}$ .
	 \choice
		 {$\mathscr{D}=\left[-1;\dfrac{5}{2}\right]$}
		{$\mathscr{D}=\mathbb{R} \setminus \{-1\}$} 
		{\True $\mathscr{D}=\mathbb{R}  \setminus \left\{-1;\dfrac{5}{2}\right\}$}
		 {$\mathscr{D}=\left(-1;\dfrac{5}{2}\right)$}
	\loigiai{ 
		Điều kiện xác định của  hàm số là $2x^2-3x-5\ne 0\Leftrightarrow \heva{&x\ne-1\\&x\ne\dfrac{5}{2}.}$\\
	Vậy tập xác định của hàm số là $\mathscr{D}=\mathbb{R} \setminus \left\{-1;\dfrac{5}{2}\right\}$.} 
\end{ex}

% nguồn: MaHoaNH_TracNghiem_Lop10
\begin{ex}%[0D3N1-2]%[ma:30aabd68]
	Tìm tập xác định  của hàm số $y=\sqrt{-8x+5}$.
	\choice 
		{$\mathscr{D}=(0;  +\infty)$}
		 {$\mathscr{D}=\mathbb{R}$}
		{$\mathscr{D}=(-\infty;0)$} 
		 {\True $\mathscr{D}=\left(-\infty;\dfrac{5}{8}\right]$}
	 \loigiai{
		Hàm  số $y=\sqrt{-8x+5}$ xác định khi $-8x+5\ge 0\Leftrightarrow x\le \dfrac{5}{8}$.\\ Vậy $\mathscr{D}=\left(-\infty;\dfrac{5}{8}\right].$
	 }
\end{ex}

% nguồn: MaHoaNH_TracNghiem_Lop10
\begin{ex}%[0D3N1-2]%[ma:b32aac76]
	Cho hệ bất phương  trình $\heva{& 2x-5y-1>0 \\& 2x+y+5>0\\&x+y+1<0}$. Trong các điểm sau, điểm nào thuộc miền nghiệm của hệ bất phương trình?
	\choice 
		 {$O\left(0;0\right)$}
		 {$M\left(1;0\right)$}
		{\True  $N\left(0;-2\right)$}
		 {$P\left(0;2\right)$}
	 \loigiai{
		Ta thay lần lượt tọa độ các điểm vào hệ bất phương  trình.\\
		Ta thấy điểm $N\left(0;-2\right)$ thỏa mãn hệ. 
	} 
\end{ex}

% nguồn: MaHoaNH_TracNghiem_Lop10
\begin{ex}%[0D3N1-2]%[ma:27f88824]
	Cho hệ bất phương trình $\heva{& x+3y-2\ge0\\ & 2x+y+1\le0 }$. Trong các điểm sau, điểm nào  thuộc miền nghiệm của hệ bất phương trình?
	\choice 
		{$M\left(0;1\right)$} 
		{\True  $N\left(-1;1\right)$}
		 {$P\left(1;3\right)$}
		{$Q\left(2;1\right)$} 
	\loigiai{ 
		Ta thay lần lượt tọa độ các điểm vào hệ  bất phương trình.\\
		Ta thấy $N\left(-1;1\right)$ thỏa  mãn.
	 }
\end{ex}

% nguồn: MaHoaNH_TracNghiem_Lop10
\begin{ex}%[0D3N1-2]%[ma:a3ab6552]
	Miền nghiệm của hệ bất phương trình  $\heva{& x>0 \\& x-y\le2\\&x+y\le1}$ chứa điểm nào sau đây?
	\choice 
		{\True $A\left(\dfrac{1}{2};-1\right)$} 
		 {$B\left(1;2\right)$}
		 {$C\left(0;2\right)$}
		 {$D\left(3;-2\right)$}
	 \loigiai{
		Thay tọa độ điểm $A\left(\dfrac{1}{2};-1\right)$  vào hệ phương trình ta thấy thỏa mãn.\\
		Suy ra điểm $A\left(\dfrac{1}{2};-1\right)$ thuộc miền nghiệm của hệ bất phương trình đã cho. 
	 }
\end{ex}

% nguồn: MaHoaNH_TracNghiem_Lop10
\begin{ex}%[0D3N1-2]%[ma:4d096e2a]
	Giá trị lớn nhất $F_{\max}$ của biểu thức $F\left(x; y\right)=x+2y$ trên miền xác định bởi hệ $\heva{&0 \leq  y \leq 4\\&x\geq0\\&x-y-1\leq 0\\&x+2y-10\leq 0} $ là
	\choice 
		{$F_{\max}=6$ } 
		{$F_{\max}=8$} 
		 {\True $F_{\max}=10$}
		 {$F_{\max}=12$}
	\loigiai{ 
		Trong mặt phẳng tọa  độ vẽ các đường thẳng $d_1:x-y-1=0$, $d_2:x+2y-10=0$, $\Delta:y=4$ \\
		Khi đó miền nghiệm của hệ bất  phương trình là phần mặt phẳng tô màu như hình vẽ.
		\begin{center}
			\begin{tikzpicture}[scale=0.6,thick,>=stealth']
					\draw[->] (-3,0) --  (12,0) node[above]{$x$};
					\foreach  \x in {1,2,3,4,5,6,7,8,9,10,11}
					 \draw[shift={(\x,0)},color=black] (0pt,2pt) -- (0pt,-2pt) node[below] {\footnotesize $\x$};
					\draw[->,color=black]  (0,-2) -- (0,8) node[right]{$y$};
					\foreach \y in {1,2,3,4,5,6,7} 
					\draw[shift={(0,\y)},color=black] (2pt,0pt) --  (-2pt,0pt) node[left] {\footnotesize $\y$};
					\node[below left] at (0,0)  {$O$};
					\clip  (-3,-2) rectangle (12,8);
					\draw[line width=1.2pt,smooth,samples=100,domain=-1:9] plot (\x,{(\x)-1}); 
					 \draw[line width=1.2pt,smooth,samples=100,domain=-1:15.3] plot (\x,{-0.5*(\x)+5});
					 \draw[line width=1.2pt,smooth,samples=100,domain=-1:15.3] plot (\x,{4});
					\fill[pattern=north east  lines] (0,0)--(0,4) --(2,4)--(4,3)--(1,0)-- cycle;
					 \node[above] at (1,0) {$A$};
					\node[right] at  (4,3) {$B$};
					 \node[above] at (2,4) {$C$};
					\node[below right]  at (0,4) {$D$};
				\end{tikzpicture}
		\end{center}
		Xét các đỉnh của miền khép kín tạo bởi hệ là $O\left(0;0\right)$, $A\left(1;0\right)$, $B\left(4;3\right)$, $C\left(2;4\right)$,  $D\left(0;4\right)$.\\
		 Ta có $\heva{&F\left(0; 0\right)=0\\ &F\left(1; 0\right)=1\\&F\left(4; 3\right)=10\\&F\left(2; 4\right)=10\\&F\left(0; 4\right)=4} $ $\Rightarrow F_{\max}=10$.
	 }
\end{ex}

% nguồn: MaHoaNH_TracNghiem_Lop10
\begin{ex}%[0D3N1-2]%[ma:852b67b7]
	Cho hệ bất  phương trình $\heva{&-2x+y\le 2 \\ & -x+2y\ge 4 \\ & x+y\ge 5}$ có miền nghiệm là miền được tô màu (miền tam giác $ABC$ có tọa độ các đỉnh là $A\left(2;0\right)$, $B\left(2;3\right)$, $C\left(1;4\right)$, bao gồm cả các cạnh) như hình vẽ:
	\begin{center}
		\begin{tikzpicture}[line join=round,line cap=round,>=stealth,font=\footnotesize,scale=0.6]
				\draw (0,0)  node [below left] {$O$};
				\draw (0,2) node  [below right] {$A$};
				 \draw (1,4) node [above] {$B$};
				\draw (2,3) node [above] {$C$}; 
				 \draw [->] (-2,0)--(4.5,0) node[below]{$x$};
				\draw[->]  (0,-1)--(0,6) node[left]{$y$};
				\clip (-2,-1)  rectangle (4.5,6);
				\foreach \x  in {1,2}
				\draw[shift={(\x,0)},color=black] (0pt,2pt) -- (0pt,-2pt) node[below] {\footnotesize  $\x$};
				\foreach \y in  {2,3,4,5}
				\draw[shift={(0,\y)},color=black] (2pt,0pt) -- (-2pt,0pt)  node[left] {\footnotesize $\y$};
				\draw[samples=100,smooth,domain=-6.5:4.5,black]  plot (\x,{2*(\x)+2});
				 \draw[samples=100,smooth,domain=-6.5:4.5,black] plot (\x,{0.5*(\x)+2});
				 \draw[samples=100,smooth,domain=-6.5:4.5,black] plot (\x,{-(\x)+5});
				\fill[pattern=north west lines,opacity=0.6] (0,2)--(2,3)--(1,4)--cycle; 
				\draw [dashed]  (1,0)--(1,4)--(0,4)
					 (2,0)--(2,3)--(0,3);
			\end{tikzpicture}
	\end{center}
	Giá trị nhỏ nhất $F_{\min}$ của biểu thức $F\left(x; y\right)$ trên  miền xác định bởi hệ trên là
	 \choice
		{\True  $F_{\min}=1$}
		 {$F_{\min}=2$}
		 {$F_{\min}=3$}
		{$F_{\min}=4$ } 
	 \loigiai{
		 Tại $A\left(2;0\right)$ ta có $F\left(2;0\right)=2$.\\
		Tại $B\left(2;3\right)$ ta có  $F\left(2;3\right)=1$.\\
		 Tại $C\left(1;4\right)$ ta có $F\left(1;4\right)=3$.\\
		 Vậy $F_{\min}=1$.
	} 
\end{ex}

% nguồn: MaHoaNH_TracNghiem_Lop10
\begin{ex}%[0D3N1-2]%[ma:5ecd12b4]
	Giá trị nhỏ nhất $F_{\min}$ của biểu thức $F\left(x; y\right)$ trên miền xác định bởi hệ $\heva{&0 \leq x \leq 10\\&0  \leq y \leq 9\\&2x+y\geq 14\\&2x+5y\geq 30} $ là
	\choice 
		 {$F_{\min}=23$}
		 {$F_{\min}=26$}
		{\True $F_{\min}=32$} 
		{$F_{\min}=67$} 
	 \loigiai{
		Trong mặt phẳng tọa độ vẽ các đường thẳng $d_1:2x+y-14=0$, $d_2:2x+5y-30=0$, $\Delta:y=9$, $\Delta'\colon  x=10$.\\
		Khi đó miền nghiệm của hệ bất phương trình là phần  mặt phẳng tô màu như hình vẽ.
		\begin{center}
			\begin{tikzpicture}[scale=0.6,thick,>=stealth']
					\draw[->] (-1,0)  -- (15.3,0) node[above]{$x$};
					\foreach \x  in {1,2,3,4,5,6,7,8,9,10,11,12,13,14,15}
					\draw[shift={(\x,0)},color=black]  (0pt,2pt) -- (0pt,-2pt) node[below] {\footnotesize $\x$};
					\draw[->,color=black]  (0,-1) -- (0,14.3) node[right]{$y$};
					\foreach \y in  {1,2,3,4,5,6,7,8,9,10,11,12,13,14}
					\draw[shift={(0,\y)},color=black] (2pt,0pt)  -- (-2pt,0pt) node[left] {\footnotesize $\y$};
					 \node[below left] at (0,0) {$O$};
					 \clip (-1,-1) rectangle (15.3,14.3);
					 \draw[line width=1.2pt,smooth,samples=100,domain=-1:9] plot (\x,{14-2*(\x)});
					 \draw[line width=1.2pt,smooth,samples=100,domain=-1:15.3] plot (\x,{6-0.4*(\x)});
					\draw[line width=1.2pt,smooth,samples=100,domain=-1:15.3] plot (\x,{9+0*(\x)}); 
					\draw[line  width=1.2pt,smooth,samples=100](10,-1)--(10,15.3);
					 \draw[line width=1.2pt,smooth,samples=100,domain=-3:9] plot (\x,{0*(\x)});
					\fill[pattern=north east lines] (5,4)--(10,2)  --(10,9)--(2.5,9)-- cycle;
					\node[below left] at (5.2,4.2)  {$A$};
					\node[right] at (10,3)  {$B$};
					\node[right] at (10,8.5)  {$C$};
					\node[left]  at (2.75,8.5) {$D$};
				\end{tikzpicture}
		\end{center}
		Xét các đỉnh của miền  khép kín tạo bởi hệ là $A\left(5; 4\right)$, $B\left(\dfrac{5}{2}; 9\right)$, $C\left(10; 9\right)$, $D\left(10; 2\right)$.\\
		Ta có $F\left(5;4\right)=32$.\\ 
		$F\left(\dfrac{5}{2}; 9\right)=37$  .\\
		$F\left(10;9\right)=67$.\\ 
		$F\left(10;2\right)=46$.\\ 
		 Vậy $F_{\min}=32$.
	} 
\end{ex}

% nguồn: MaHoaNH_TracNghiem_Lop10
\begin{ex}%[0D3H1-2]%[ma:aa8e30ca]
	Tập xác định của hàm số $y=\dfrac{1}{\sqrt {x-3}}+\sqrt  {x-1}$ là
	\choice 
		{\True $\mathscr{D}=(3;+\infty)$} 
		 {$\mathscr{D}=[1;+\infty)$}
		 {$\mathscr{D}=(1;+\infty)\setminus\{3\}$}
		{$\mathscr{D}=[1;+\infty)\setminus\{3\}$} 
	\loigiai{Điều kiện $\heva{&x-3>0\\&x-1\geq 0}\Leftrightarrow\heva{&x>3\\&x\geq 1}\Leftrightarrow  x>3$.}
\end{ex}

% nguồn: MaHoaNH_TracNghiem_Lop10
\begin{ex}%[0D3H1-2]%[ma:b91b65e2]
	Tập xác định của hàm số $y=\dfrac{5}{x^2-4}$ là 
	\choice 
		{$\mathbb{R} \setminus  \{-2\}$}
		 {\True $\mathbb{R} \setminus \{-2;2\}$}
		{$\mathbb{R}  \setminus \{2\}$}
		{$\mathbb{R}$} 
	 \loigiai{
		Hàm số đã cho  xác định khi $x^2-4 \neq 0 \Leftrightarrow \heva{&x \neq 2 \\ &x \neq -2}$ \\
	Vậy tập xác  định của hàm số là $\mathscr{D}=\mathbb{R} \setminus \{-2;2\}$.}
\end{ex}

% nguồn: MaHoaNH_TracNghiem_Lop10
\begin{ex}%[0D3H1-2]%[ma:fd8ad5a0]
	Tìm tập  xác định của hàm số $y=\dfrac{x}{x^2-2023^2}$.
	 \choice
		 {\True $\mathscr{D}=\mathbb{R}\backslash\big\{\pm 2023\big\}$}
		 {$\mathscr{D}=\mathbb{R}$}
		 {$\mathscr{D}=\big\{\pm 2023\big\}$}
		 {$\mathscr{D}=\mathbb{R}\backslash\big\{2023\big\}$}
	 \loigiai{
		Hàm số có nghĩa khi và chỉ khi $x^2-2023^2\neq 0\Leftrightarrow x \neq \pm  2023$.
	 }
\end{ex}

% nguồn: MaHoaNH_TracNghiem_Lop10
\begin{ex}%[0D3H1-2]%[ma:16000e76]
	Tìm tất cả các  giá trị của tham số $m$ để hàm số $y=\dfrac{x+1}{x-2m+1}$ xác định trên $[0;1)$.
	\choice 
		{$m <\dfrac{1}{2}$} 
		 {$m\ge 1$}
		{$m\ge 2$  hoặc $m < 1$}
		 {\True $m <\dfrac{1}{2}$ hoặc $m\ge 1$}
	 \loigiai{
		Tập  xác định của hàm số $\mathscr{D}=\mathbb{R}\setminus\{{2m-1}\}$.\\
		Hàm số xác định trên  $[0;1)$ khi và chỉ khi
	$$2m-1\notin[0;1) \Leftrightarrow \hoac{&2m-1 < 0\\ &2m-1\ge 1} \Leftrightarrow \hoac{&m <\dfrac{1}{2}\\ &m\ge 1.}$$} 
\end{ex}

% nguồn: MaHoaNH_TracNghiem_Lop10
\begin{ex}%[0D3V1-2]%[ma:f79c18cc]
	Tìm tất cả các giá trị thực của tham số $m$ để hàm số $y=\dfrac{x^2+2m+2}{x-m}$ xác định  trên khoảng $(-1;0)$.
	\choice 
		{$\hoac{&m > 0\\ &m <-1}$} 
		{$m\le-1$} 
		{\True  $\hoac{&m\ge 0\\ &m\le-1}$}
		 {$m\ge 0$}
	 \loigiai{
		 Hàm số đã cho xác định khi và chỉ khi $x\ne m$.\\
		Khi đó tập xác định của hàm  số là $\mathscr{D}=(-\infty;m)\cup(m;+\infty)$.\\
	Yêu cầu  bài toán $\Leftrightarrow(-1;0)\subset D \Leftrightarrow \hoac{&m\ge 0\\ &m\le-1.}$}
\end{ex}

% nguồn: MaHoaNH_TracNghiem_Lop10
\begin{ex}%[0D3B1-2]%[ma:f50d3c23]
	Không gian mẫu của phép thử ``\textit{gieo một con súc sắc $6$  mặt hai lần liên tiếp}'' có bao nhiêu phần tử?
	 \choice
		 {$12$}
		{\True  $36$}
		{$24$} 
		 {$30$}
	 \loigiai{
		Phép thử ``\textit{gieo một con súc sắc $6$ mặt  hai lần liên tiếp}'' có $6^2=36$ phần tử.
	} 
\end{ex}

% nguồn: MaHoaNH_TracNghiem_Lop10
\begin{ex}%[0D3Y1-2]%[ma:519454e6]
	Gieo con súc sắc $2$ lần. Biến cố nào dưới  đây là biến cố tổng số chấm trong hai lần gieo bằng $7$?
	\choice 
		 {\True $A=\left\{(1;6),(6;1),(3;4),(4;3),(5;2),(2;5)\right\}$}
		{$B=\left\{(1;6),(3;4),(5;2)\right\}$} 
		 {$C=\left\{(1;6),(6;1),(3;4),(4;3)\right\}$}
		{$D=\left\{(2;6),(6;2),(3;4),(4;3),(5;2),(2;5)\right\}$} 
	\loigiai{ 
		Ta  có $A=\left\{(1;6),(6;1),(3;4),(4;3),(5;2),(2;5)\right\}$.
	} 
\end{ex}

% nguồn: MaHoaNH_TracNghiem_Lop10
\begin{ex}%[0D3Y1-2]%[ma:c98b108d]
	Tập xác định của hàm số $y=\sqrt{x-3}$  là
	 \choice
		{$(3;+\infty)$} 
		 {$(-\infty;3)$}
		 {\True $[3;+\infty)$}
		 {$(-\infty;3]$}
	 \loigiai{
		Điều kiện xác định $x-3 \ge 0 \Leftrightarrow x \ge 3$. Vậy tập xác định là $\mathscr{D}=[3;+\infty)$. 
	} 
\end{ex}

% nguồn: MaHoaNH_TracNghiem_Lop10
\begin{ex}%[0D3Y1-2]%[ma:461f34c2]
	Gieo một con xúc xắc cân đối và đồng chất ba lần. Số  phần tử của không gian mẫu là
	 \choice
		 {\True $6^3$}
		 {$6^2$}
		{$6$} 
		 {$6^4$}
	 \loigiai{
		Ta có $n\left(\Omega\right)=6\cdot 6\cdot 6=6^3$. 
	 }
\end{ex}

% nguồn: MaHoaNH_TracNghiem_Lop10
\begin{ex}%[0D3Y1-2]%[ma:d9837813]
	Tập xác định  của hàm số $y=\dfrac{3}{x-5}$ là
	 \choice
		 {$D=(5;+\infty)$}
		 {$D= \mathbb{R}$}
		{\True  $D= \mathbb{R} \backslash\{5\}$}
		{$D= \mathbb{R}  \backslash\{-5\}$}
	 \loigiai{
		Điều kiện xác định $x-5\neq 0\Leftrightarrow x\neq 5$.\\ 
		Tập xác định của hàm số $y=\dfrac{3}{x-5}$  là $D= \mathbb{R} \backslash\{5\}$.
	} 
\end{ex}

% nguồn: MaHoaNH_TracNghiem_Lop10
\begin{ex}%[0D3B1-2]%[ma:f83f5624]
	Rút ngẫu nhiên một hộp chứa $9$ tấm thẻ được đánh số $1$; $2$; $3$; $4$; $5$; $6$; $7$; $8$; $9$. Gọi $A$ là biến cố: “Rút được thẻ  ghi số chia hết cho $3$”. Biến cố đối của biến cố $A$ là
	\choice 
		 {$\overline{A}=\left\{1;2;4;5;7;8\right\}$}
		 {$\overline{A}=\left\{3;6;9\right\}$}
		{\True  $\overline{A}=\left\{1;2;4;5;7;8;10\right\}$}
		 {$\overline{A}=\left\{2;4;6;8\right\}$}
	 \loigiai{
		Biến cố đối của biến cố $A$ là “Rút được thẻ ghi số không chia  hết cho $3$”. \\
		Do  đó, $\overline{A}=\left\{1;2;4;5;7;8;10\right\}$.
	 }
\end{ex}

% nguồn: MaHoaNH_TracNghiem_Lop10
\begin{ex}%[0D3B1-2]%[ma:622468d6]
	Tìm tập xác định $\mathscr{D}$ của hàm số $y=\sqrt{x-1}$. 
	\choice 
		 {$\mathscr{D}=\mathbb{R}\setminus\{1\}$}
		 {$\mathscr{D}=(1;+\infty)$}
		{\True  $\mathscr{D}=[1;+\infty)$}
		{$\mathscr{D}=(-\infty;1]$} 
	\loigiai{ 
		Hàm  số xác định $\Leftrightarrow x-1\geq 0\Leftrightarrow x\geq 1$.\\
		Vậy tập xác định là $\mathscr{D}=[1;+\infty)$. 
	 }
\end{ex}

% nguồn: MaHoaNH_TracNghiem_Lop10
\begin{ex}%[0D3Y1-2]%[ma:a1098346]
	Lớp 10A có $40$ học sinh với $24$ nữ và $16$  nam. Giáo viên chủ nhiệm cần chọn ra $4$ học sinh. Số phần tử của biến cố $A\colon$ "trong số 4 học sinh được chọn có ít nhất $2$ nam" là
	\choice 
		{$48830$} 
		{\True $48380$} 
		 {$46560$}
		 {$33120$}
	 \loigiai{
		Số phần tử của biến  cố $A$ là
		$$n(A)=\mathrm{C}_{40}^4-(\mathrm{C}_{24}^4+\mathrm{C}_{24}^3\cdot\mathrm{C}_{16}^1)=48380.$$ 
	} 
\end{ex}

% nguồn: MaHoaNH_TracNghiem_Lop10
\begin{ex}%[0D3Y1-2]%[ma:d8ee67d2]
	Tìm tập xác định của  hàm số $y=\dfrac{x}{x^2-2023^2}$ là
	 \choice
		{\True  $\mathcal{D}=\mathbb{R}\backslash\{-2023; 2023\}$}
		{$\mathcal{D}=\mathbb{R}$} 
		{$\mathcal{D}=\{-2023;  2023\}$}
		 {$\mathcal{D}=\mathbb{R} \backslash\{2023\}$}
	 \loigiai{
		Hàm  số xác định khi $x^2-2023^2 \neq 0 \Leftrightarrow x \neq \pm 2023$.\\
		Vậy tập xác định  của hàm số là $\mathcal{D}=\mathbb{R}\backslash\{-2023; 2023\}$.
	} 
\end{ex}

% nguồn: MaHoaNH_TracNghiem_Lop10
\begin{ex}%[0D3B1-2]%[ma:c0cf668e]
	Tập xác định $\mathscr{D}$  của hàm số $y=\dfrac{3x-1}{2x-2}$ là
	\choice 
		{\True  $\mathscr{D}=\mathbb{R}\setminus\{1\}$}
		 {$\mathscr{D}=[1;+\infty)$}
		{$\mathscr{D}=(1;+\infty)$} 
		{$\mathscr{D}=\mathbb{R}$} 
	\loigiai{ 
		Điều  kiện $2x-2\neq 0\Leftrightarrow x\neq 1 $.\\
		Vậy tập xác định là  $\mathscr{D}=\mathbb{R}\setminus\{1\}$.
	 }
\end{ex}

% nguồn: MaHoaNH_TracNghiem_Lop10
\begin{ex}%[0D3B1-2]%[ma:e730464e]
	Tập  xác định của hàm số $f(x)=\sqrt{3-x}+\dfrac{1}{\sqrt{x-1}}$ là
	\choice 
		{\True  $\mathscr{D}=(1;3] $}
		{$\mathscr{D}=(-\infty;1)\cup [3;+\infty)$} 
		 {$\mathscr{D}=[1;3]$}
		 {$\mathscr{D}=\varnothing$}
	 \loigiai{
		Điều  kiện $\heva{&3-x\geq 0\\&x-1>0}\Leftrightarrow\heva{&x\leq 3\\&x>1.}$ \\
		 Vậy $\mathscr{D}=(1;3] $.
	} 
\end{ex}

% nguồn: MaHoaNH_TracNghiem_Lop10
\begin{ex}%[0D3Y1-2]%[ma:fffbf3d6]
	Tập xác định $\mathscr{D}$  của hàm số $y=\sqrt{3x-1}$ là
	\choice 
		{$\mathscr{D}=(0;+\infty)$} 
		{$\mathscr{D}=[0;+\infty)$} 
		 {$\mathscr{D}=\left(\dfrac{1}{3};+\infty\right)$}
		 {\True $\mathscr{D}=\left[\dfrac{1}{3};+\infty\right)$}
	 \loigiai{
		Điều kiện xác định của hàm số $y=\sqrt{3x-1}$ là $3x-1 \geq 0 \Leftrightarrow x \geq \dfrac{1}{3}$.\\ 
		Tập xác  định của hàm số là $\mathscr{D}=\left[\dfrac{1}{3};+\infty\right)$.
	 }
\end{ex}

% nguồn: MaHoaNH_TracNghiem_Lop10
\begin{ex}%[0D3K1-2]%[ma:79e963f7]
	Tìm các giá trị thực  của tham số $m$ để hàm số $y=\dfrac{x+m+2}{x-m}$ xác định trên $(-1;3)$.
	 \def\dotEX{}
	 \choice
		 {$-1<m<3$.}
		{$\hoac{&m<-1\\&m>3.}$} 
		{\True $\hoac{&m \leq  -1\\&m \geq 3.}$}
		{$\heva{&m \leq -1\\&m  \geq 3.}$}
	 \loigiai{
		Điều kiện xác định của $y=\dfrac{x+m+2}{x-m}$  là $x \neq m$.\\
		Do đó hàm số xác định trên khoảng $(-1;3)$ khi và chỉ khi  $m \notin (-1;3) \Leftrightarrow \hoac{&m \leq -1\\&m \geq 3.}$
	} 
\end{ex}

% nguồn: MaHoaNH_TracNghiem_Lop10
\begin{ex}%[0D3Y1-2]%[ma:dc310fc3]
	Hàm số  $y=\dfrac{2}{2-x} + \sqrt{4-x^2}$ có tập xác định là
	\choice 
		 {$(-2;2]$}
		{$(-2;2)$} 
		{$[-2;2]$} 
		 {\True $[-2;2)$}
	\loigiai{ 
		Điều kiện xác  định:
		$\heva{&2-x \ne 0  \\& 4-x^2 \geq 0} \Leftrightarrow \heva{&x \ne 2 \\& -2 \leq x \leq 2} \Leftrightarrow -2 \leq x <2$.\\
		Vậy tập xác định của hàm  số đã cho là $[-2;2)$.
	} 
\end{ex}

% nguồn: MaHoaNH_TracNghiem_Lop10
\begin{ex}%[0D3B1-2]%[ma:66a83660]
	Trong các hàm số sau, hàm số nào có  tập xác định là $\mathbb{R}$?
	\choice 
		{$y=\dfrac{x+2}{x-1}$} 
		{$y=\dfrac{x^2+2}{x}$} 
		 {$y=\dfrac{2x+3}{x^2}$}
		{\True $y=x^3+3x^2-1$} 
	 \loigiai{
		Hàm  số $y=\dfrac{x+2}{x-1}$ có tập xác định $\mathbb{R} \setminus \{1\}$.\\
		Hàm số $y=\dfrac{x^2+2}{x}$ có tập xác định $\mathbb{R} \setminus \{0\}$.\\ 
		Hàm số $y=\dfrac{2x+3}{x^2}$  có tập xác định $\mathbb{R} \setminus \{0\}$.\\
		Hàm số $y=x^3+3x^2-1$ có tập xác  định $\mathbb{R}$.
	} 
\end{ex}

% nguồn: MaHoaNH_TracNghiem_Lop10
\begin{ex}%[0D3G1-2]%[ma:dd62db59]
	Có tât cả bao nhiêu giá trị nguyên của tham số $m$ để hàm số $y=\dfrac{2x+1}{\sqrt{(m+2)x^2-2(m+2)x+4}}$ xác  định với mọi $x\in\mathbb{R}$ .
	\choice 
		 {$3$}
		{\True  $4$}
		{$5$} 
		{$6$} 
	\loigiai{ 
		Yêu cầu bài toán $\Leftrightarrow f(x)=(m+2)x^2-2(m+2)x+4>0,\forall x\in\mathbb{R}$.  $\qquad(1)$
		\begin{itemize}
			\item Trường  hợp $1$: $m+2=0\Leftrightarrow m=-2$.\\
			Ta có $(1)\Leftrightarrow 4>0,\forall  x\in\mathbb{R}$.\\
			 Nhận $m=-2$.
			\item Trường hợp $2$: $m+2\neq 0\Leftrightarrow m\neq -2$.\\ 
			Ta có  $(1)\Leftrightarrow \heva{&m+2>0\\&(m+2)^2-4(m+2)<0}\Leftrightarrow\heva{&m>-2\\&0<m+2<4}\Leftrightarrow -2<m<2$.\\
			 Mà $m$ nguyên suy ra $m\in\{-1;0;1\}$.
		\end{itemize}
		Vậy có $4$ giá trị $m\in\{-2;-1;0;1\}$. 
	 }
\end{ex}


%%%%%%%%%%%%%%%% ĐÚNG/SAI %%%%%%%%%%%%%%%%

% nguồn: MaHoaNH_DungSai_Lop10
\begin{ex}%[0D3N1-2]%[ma:14af49d3]
	Xét tính đúng, sai của các khẳng định sau 
	 \choiceTF
		{\True  Tập xác định của hàm số $y=\dfrac{x^2-2x+3}{x+7}$ là $\mathscr{D}=\mathbb{R}\setminus\{-7\}$}
		{Tập xác định của hàm  số $y=\sqrt{x-6}$ là $\mathscr{D}=\mathbb{R}$}
		 {\True Tập xác định của hàm số $y=3x^2$ là $\mathscr{D}=\mathbb{R}$}
		{Tập xác định  của hàm số $y=\sqrt{1-x^2}$ là $\mathscr{D}=(-1;1)$}
	\loigiai{ 
		\begin{itemchoice}
			\itemch Đúng. Hàm số $y=\dfrac{x^2-2x+3}{x+7}$ xác định khi $x+7\ne  0 \Leftrightarrow x\ne-7$.\\
			 Vậy tập xác định của hàm số $y=\dfrac{x^2-2x+3}{x+7}$ là $\mathscr{D}=\mathbb{R}\setminus\{-7\}$.
			\itemch Sai. Hàm số $y=\sqrt{x-6}$ xác định khi $x-6\ge 0\Leftrightarrow x\ge 6$.\\ 
			Vậy tập xác định  của hàm số $y=\sqrt{x-6}$ là $\mathscr{D}=[6;+\infty)$.
			\itemch Đúng. Tập xác định của hàm số  $y=3x^2$ là $\mathscr{D}=\mathbb{R}$.
			\itemch Sai. Hàm số  $y=\sqrt{1-x^2}$ xác định khi $1-x^2\ge 0\Leftrightarrow-1\le x\le 1$.\\
			Vậy tập xác định  của hàm số $y=\sqrt{1-x^2}$ là $\mathscr{D}=[-1;1]$.
		\end{itemchoice}
	 }
\end{ex}

% nguồn: MaHoaNH_DungSai_Lop10
\begin{ex}%[0D3H1-2]%[ma:d6c58884]
	Cho hàm số $f(x)=\sqrt{1-x}+\sqrt{x+3}$. Mỗi kết quả dưới đây đúng  hay sai?
	\choiceTF 
		{$f(1)=3$} 
		{\True $f(2x)=\sqrt{1-2x}+\sqrt{2x+3}$} 
		{\True  Tập xác định $\mathscr{D}=\left[-3;1\right]$}
		{\True  $f(x)\ge 2$, $\forall x\in\mathscr{D}$}
	\loigiai{ 
		\begin{itemchoice}
			\itemch  $f(1)=\sqrt{1-1}+\sqrt{1+3}=2$.
			\itemch  $f(2x)=\sqrt{1-2x}+\sqrt{2x+3}$.
			\itemch Điều kiện  $\heva{& 1- x\ge 0\\& x+3 \ge 0}$ nên $\mathscr{D}=\left[-3;1\right]$
			\itemch $\left(f(x)\right)^2=4+2\sqrt{\left(1-x\right)\left(x+3\right)}\ge 4\Rightarrow  f(x)\ge 2$ (vì $f(x)>0$).
		\end{itemchoice}
	} 
\end{ex}

% nguồn: MaHoaNH_DungSai_Lop10
\begin{ex}%[0D3H1-2]%[ma:f4bb73d2]
	Cho hàm số $f(x)=\heva{&2x+1,&\text{khi} \, x\ge  -1\\& -2-x, & \text{khi} \, x<-1}$. Mỗi kết quả dưới đây đúng hay sai?
	 \choiceTF
		 {\True $f(1)=3$}
		 {$f(2)+f(-2)=0$}
		{Tập xác định  $\mathscr{D}=\mathbb{R}\setminus\left\{1\right\}$}
		{\True $f(x^2)=2x^2+1$  khi $x\ge -1$}
	 \loigiai{
		\begin{itemchoice}
			\itemch $f(1)=2\cdot 1+1 =  3$.
			 \itemch $f(2)+f(-2)=2\cdot 2+1+(-2)-(-2)=5$.
			\itemch  $\mathscr{D}=\left(-\infty;-1\right)\cup\left[-1;+\infty\right)=\mathbb{R}$
			\itemch $f(x^2)=2\cdot  x^2+1$.
		\end{itemchoice}
	} 
\end{ex}

% nguồn: MaHoaNH_DungSai_Lop10
\begin{ex}%[0D3H1-2]%[ma:7684ac3c]
	Cho hàm số $f(x)=x^3+1$. Mỗi kết quả dưới  đây đúng hay sai?
	\choiceTF 
		{Phương trình $f(x)=0$  vô nghiệm}
		 {\True Tập xác định $\mathscr{D}=\mathbb{R}$}
		{\True Hàm số đồng  biến trên $\left(-\infty;+\infty\right)$}
		{$f(x)+1\ge 1$,  $\forall x\in\mathscr{D}$}
	\loigiai{ 
		\begin{itemchoice}
			 \itemch $f(x)=0\Leftrightarrow x^3+1=0\Leftrightarrow x=-1$.
			\itemch $\mathscr{D}=\mathbb{R}$. 
			\itemch Với  $x_1$, $x_2\in\mathbb{R}$ sao cho $x_1\ne x_2$.\\
			Xét $M=\dfrac{f(x_1)-f(x_2)}{x_1-x_2}=\dfrac{x_1^3-x_2^3}{x_1-x_2}=x_1^2+x_1x_2+x_2^2>0$.\\ 
			Hàm số  đồng biến trên $\left(-\infty;+\infty\right)$.
			\itemch Ta có  $f(-2)=-7$.
		\end{itemchoice}
	 }
\end{ex}

% nguồn: MaHoaNH_DungSai_Lop10
\begin{ex}%[0D3H1-2]%[ma:0c9211f1]
	Cho hàm số $f(x)=\sqrt{1-x}$. Mỗi kết quả dưới đây đúng hay sai? 
	 \choiceTF
		{\True $f(0)=1$} 
		{\True Tập  xác định $\mathscr{D}=\left(-\infty;1\right]$}
		{Hàm số đồng biến trên  $\left(-\infty;1\right)$}
		{\True  $f(x^2)\le 1$, $\forall x\in\mathscr{D}$}
	 \loigiai{
		\begin{itemchoice}
			\itemch $f(0)=\sqrt{1-0}=1$. 
			\itemch Điều  kiện $1-x\ge 0\Leftrightarrow x\le 1\Rightarrow \mathscr{D}=\left(-\infty;1\right]$.
			\itemch Với mọi $x_1$,  $x_2\in\mathscr{D}$ sao cho $x_1\ne x_2$.\\
			Xét  $M=\dfrac{f(x_1)-f(x_2)}{x_1-x_2}=\dfrac{\sqrt{1-x_1}-\sqrt{1-x_2}}{x_1-x_2}=\dfrac{-1}{\sqrt{1-x_1}+\sqrt{1-x_2}}<0$.\\
			Vậy hàm số nghịch  biến trên $\left(-\infty;1\right)$.
			\itemch $f(x^2)=\sqrt{1-x^2}\le 1$. 
		\end{itemchoice}
	} 
\end{ex}

% nguồn: MaHoaNH_DungSai_Lop10
\begin{ex}%[0D3H1-2]%[ma:8f92511d]
	Cho hàm số  $f(x)=\dfrac{x-1}{x+1}$. Mỗi kết quả dưới đây đúng hay sai?
	 \choiceTF
		{\True $f(0)=-1$} 
		{Tập xác  định $\mathscr{D}=\left(-\infty;-1\right)$}
		{ Hàm số đồng biến trên  $\left(-\infty;1\right)$}
		{\True Tồn tại $4$ số nguyên $x$ để $f(x)$ là  số nguyên}
	 \loigiai{
		\begin{itemchoice}
			\itemch $f(0)=\dfrac{0-1}{0+1}=-1$. 
			\itemch Điều kiện $x+1\ne 0\Leftrightarrow  x\ne -1$.\\
			Vậy $\mathscr{D}=\left(-\infty;-1\right)\cup\left(-1;+\infty\right)$. 
			\itemch Với  $x_1$, $x_2\in\mathscr{D}$ sao cho $x_1\ne x_2$.\\
			 Xét biểu thức $M=\dfrac{f(x_1)-f(x_2)}{x_1-x_2}=\dfrac{\dfrac{x_1-1}{x_1+1}-\dfrac{x_2-1}{x_2+1}}{x_1-x_2}=\dfrac{2\left(x_1-x_2\right)}{\left(x_1+1\right)\left(x_2+1\right)\left(x_1-x_2\right)}=\dfrac{2}{\left(x_1+1\right)\left(x_2+1\right)}$.\\
			Với $x_1$, $x_2\in\left(-\infty;-1\right)$ thì $\left(x_1+1\right)\left(x_2+1\right)>0$ nên  hàm số đồng biến $\left(-\infty;-1\right)$.\\
			Với  $x_1$, $x_2\in\left(-1;+\infty\right)$ thì $\left(x_1+1\right)\left(x_2+1\right)>0$ nên hàm số đồng biến $\left(-1;+\infty\right)$.\\
			\itemch $f(x)=\dfrac{x+1-2}{x+1}=1-\dfrac{2}{x+1}\in\mathbb{Z}$ khi $x+1$ là ước nguyên của  $2$.\\
			Vậy tồn  tại $4$ số nguyên $x$ để $f(x)$ là số nguyên.
		\end{itemchoice}
	} 
\end{ex}

% nguồn: MaHoaNH_DungSai_Lop10
\begin{ex}%[0D3H1-2]%[ma:5fbe1f3c]
	Xét  tính đúng, sai của các khẳng định sau:
	\choiceTF 
		{\True Hàm  số $y=-x^3-|x|+1$ có tập xác định $\mathscr{D}=\mathbb{R}$}
		{Hàm số $y=\dfrac{3x-2}{x^2-x}$ có tập xác định  $\mathscr{D}=\mathbb{R} \setminus \{1\}$}
		{\True  Hàm số $y=\sqrt{x+2}+\sqrt{x-1}$ có tập xác định $\mathscr{D}=[1;+\infty)$}
		{Hàm số $y=\dfrac{1}{\sqrt{2-|x|}}$ có tập xác  định $\mathscr{D}=[-2;2]$}
	 \loigiai{
		\begin{itemchoice}
			\itemch Đúng. Hàm  số xác định với mọi $x$ thuộc $\mathbb{R}$.\\
			Tập xác định hàm số là  $\mathscr{D}=\mathbb{R}$.
			\itemch Sai. Hàm số xác định $\Leftrightarrow  x^2-x\ne 0
			\Leftrightarrow \heva{&x\ne 0\\ &x\ne  1.}$\\
			Tập xác  định hàm số là $\mathscr{D}=\mathbb{R} \setminus \{0; 1\}$.
			\itemch Đúng. Hàm số xác định $\Leftrightarrow \heva{&x+2\ge 0\\& x-1\ge 0} 
			\Leftrightarrow \heva{&x\ge-2\\ &x\ge 1} \Leftrightarrow x\ge  1$.\\
			Tập xác  định hàm số là $\mathscr{D}=[1;+\infty)$.
			\itemch Sai. Hàm số  xác định $\Leftrightarrow 2-|x| > 0 \Leftrightarrow |x| < 2 \Leftrightarrow -2<x<2$.\\
			 Tập xác định hàm số là $\mathscr{D}=(-2;2)$.
	\end{itemchoice}}
\end{ex}

% nguồn: MaHoaNH_DungSai_Lop10
\begin{ex}%[0D3V1-2]%[ma:46430bba]
	Cho hàm số $f(x)=x^2-2mx+4$. Mỗi kết quả dưới đây đúng hay sai? 
	 \choiceTF
		{\True Có $3$ số nguyên của tham  số $m$ để phương trình $f(x)=0$ vô nghiệm}
		 {\True Tập xác định $\mathscr{D}=\mathbb{R}$}
		{Hàm số đồng biến trên  $\left(-\infty;-1\right)$}
		{\True  Tồn tại $4$ giá trị nguyên dương của tham số $m$ để $f(x)$ đồng biến trên $\left(4;+\infty\right)$}
	 \loigiai{
		\begin{itemchoice}
			\itemch $f(x)=0\Leftrightarrow x^2-2mx+4=0$. Để phương trình vô nghiệm thì $\Delta=m^2-4<0\Leftrightarrow -2<m<2$. 
			\itemch $\mathscr{D}=\mathbb{R}$. 
			 \itemch Hàm số nghịch biến $\left(-\infty;m\right)$.
			\itemch Hàm số đồng biến  trên $\left(m;+\infty\right)$ nên suy ra $m\le 4$.
		\end{itemchoice}
	 }
\end{ex}


%%%%%%%%%%%%%%%% TRẢ LỜI NGẮN %%%%%%%%%%%%%%%%

% nguồn: MaHoaNH_TraLoiNgan_Lop10
\begin{bt}%[0D3H1-2]%[ma:a5866526]
	Có bao nhiêu số nguyên $m$ thuộc đoạn $[1;8]$ để hàm  số $y=\sqrt{3x-m}$, (với $m$ là tham số) xác định trên tập $(1;+\infty)$.
	\shortans{3} 
	 \loigiai{
		Điều kiện  xác định của hàm số là $3x-m\ge 0 \Leftrightarrow x\ge\dfrac{m}{3}$.\\
		Suy ra $\mathscr{D}=\left[\dfrac{m}{3};+\infty\right)$.\\ 
		Để  hàm số xác định trên $(1;+\infty)$ thì $$(1;+\infty)\subset\left[\dfrac{m}{3};+\infty\right) \Leftrightarrow\dfrac{m}{3}\le 1 \Leftrightarrow m\le 3.$$
		Với $m$ thuộc đoạn $[1;8]$  nên ta có $m=\{1;2;3\}$.\\
	 Vậy có $3$ số nguyên $m$ thỏa mãn yêu cầu.}
\end{bt}

% nguồn: MaHoaNH_TraLoiNgan_Lop10
\begin{bt}%[0D3H1-2]%[ma:56720b65]
	 Tập xác định của hàm số $y=\sqrt{\dfrac{3-x}{x+2}}+\sqrt[3]{x-1}$ có dạng $(a;b]$. Khi đó $a+b$ bằng
	\shortans{1} 
	\loigiai{ 
		Hàm số xác định khi và  chỉ khi $\dfrac{3-x}{x+2}\ge 0 \Leftrightarrow
		 \hoac{
		\heva{&3-x\ge 0\\ &x+2  > 0}\\
		\heva{&3-x\le 0\\ &x+2 <  0}}
		 \Leftrightarrow \hoac{
		&\heva{&x\le 3\\ &x  >-2}\\
		&\heva{&x\ge 3\\ &x <-2}} 
		\Leftrightarrow-2 < x\le  3$.\\
		Tập xác định  hàm số là $\mathscr{D}=(-2;3]$.\\
	Vậy  $a+b=-2+3=1$.}
\end{bt}

% nguồn: MaHoaNH_TraLoiNgan_Lop10
\begin{bt}%[0D3V1-2]%[ma:ada63961]
	Có bao nhiêu giá trị nguyên của tham số $m$ thuộc khoảng $(-2;5)$  để hàm số \break $y=\dfrac{x+2m+2}{x-m}$ xác định trên $(-1;0)$.
	 \shortans{6}
	 \loigiai{
		 Điều kiện xác định của hàm số là $x\ne m$.\\
		Để hàm số xác định trên $(-1;0)  \Leftrightarrow m\notin (-1;0)
		\Leftrightarrow\hoac{&m\ge 0\\ &m\le-1.}$\\ 
		Với $m$ thuộc khoảng $(-2;5)$ ta có $m=\{-1;0;1;2;3;4\}$.\\ 
	Vậy có $6$ giá trị  của $m$.}
\end{bt}

% nguồn: MaHoaNH_TraLoiNgan_Lop10
\begin{bt}%[0D3V1-2]%[ma:7b727fb1]
	Có bao nhiêu  giá trị nguyên âm của tham số $m$ để hàm số $y=\sqrt{x^2-2m x-2m+3}$ có tập xác định là $\mathbb{R}$.
	\shortans{3} 
	\loigiai{ 
		Hàm số $y=\sqrt{x^2-2m x-2m+3}$ có tập xác định là  $\mathbb{R}$ khi $x^2-2m x-2m+3\geq 0$ với mọi $x \in \mathbb{R}$. Điều này tương đương với
		$$\heva{&\Delta' \leq 0\\&a > 0} \Leftrightarrow\heva{&m^2+2m-3\leq  0 \\&1> 0} \Leftrightarrow -3\leq m \leq 1.$$
		Do  $m$ nguyên âm nên $m \in\{-3;-2;-1\}$.
		Vậy có $3$ giá trị nguyên âm  của $m$ thỏa yêu cầu bài toán.
	 }
\end{bt}

% nguồn: MaHoaNH_TraLoiNgan_Lop10
\begin{bt}%[0D3V1-2]%[ma:f18012cf]
	 Cho hàm số $y=\sqrt{2x-3m+4}+\dfrac{x}{x+m-1}$ với $m$ là tham số. Tìm tổng các giá trị của $m$ để hàm số có tập xác định là $[0;+\infty)$ (kết quả dưới dạng số thập phân làm tròn đến hàng phần trăm).
	 \shortans{1,33}
	 \loigiai{
		Hàm số xác định khi  và chỉ khi $\heva{&2x-3m+4\ge 0\\ &x+m-1\ne 0} \Leftrightarrow \heva{&x\ge\dfrac{3m-4}{2}\\ &x\ne 1-m.}$
		\begin{itemize}
			 \item \textbf{Trường hợp 1:} $1-m\ge\dfrac{3m-4}{2}\Leftrightarrow m\le\dfrac{6}{5}$, khi đó tập xác định của hàm số là $\mathscr{D}=\left[\dfrac{3m-4}{2};+\infty\right)\setminus\{ 1-m\}$.\\
			Tập xác định này không thể bằng  $[0;+\infty)$ theo đề bài.\\
			Do đó $m\le\dfrac{6}{5}$ không thỏa  mãn.
			\item \textbf{Trường hợp  2:} $1-m <\dfrac{3m-4}{2}\Leftrightarrow m >\dfrac{6}{5}$, khi đó tập xác định của hàm số là $\mathscr{D}=\left[\dfrac{3m-4}{2};+\infty\right)$.\\
			Do đó, hàm số có tập xác định là $[0;+\infty)\Leftrightarrow\dfrac{3m-4}{2}=0\Leftrightarrow m=\dfrac{4}{3}$ (thỏa  mãn).\\
			Vậy $m=\dfrac{4}{3}\approx 1{,}33$ là giá trị  cần tìm.
		\end{itemize}
	} 
\end{bt}
